% Grundlegende Dokumenteneigenschaften gemäß DHBW-Vorgaben
\documentclass[a4paper,fontsize=11pt,oneside,parskip=half,headings=normal]{scrreprt} 
% \usepackage{showframe} % nur für Kontrolle der Ränder 

%%% Präambel einbinden (mit Festlegungen gemäß DHBW-Vorgaben) %%%
%%% Präambel %%%
% hier sollten keine Änderungen erforderlich sein
%
\usepackage[utf8]{inputenc}   % Zeichencodierung UTF-8 für Eingabe-Dateien
\usepackage[T1]{fontenc}      % Darstellung von Umlauten im PDF

\usepackage{listings}         % für Einbindung von Code-Listings
\lstset{numbers=left,numberstyle=\tiny,numbersep=5pt,texcl=true}
\lstset{literate=             % erlaubt Sonderzeichen in Code-Listings 
  {Ö}{{\"O}}1
{Ä}{{\"A}}1
{Ü}{{\"U}}1
{ß}{{\ss}}2
{ü}{{\"u}}1
{ä}{{\"a}}1
{ö}{{\"o}}1
{€}{{\euro}}1
}

\usepackage[
  inner=35mm,outer=15mm,top=25mm,
  bottom=20mm,foot=12mm,includefoot
]{geometry}                 % Einstellungen für Ränder

\usepackage[ngerman]{babel} % Spracheinstellungen Deutsch
\usepackage[babel,german=quotes]{csquotes} % deutsche Anf.zeichen
\usepackage{enumerate}      % anpassbare Nummerier./Aufz.
\usepackage{graphicx}       % Einbinden von Grafiken
\usepackage[onehalfspacing]{setspace} % anderthalbzeilig

\usepackage{blindtext}      % Textgenerierung für Testzwecke
\usepackage{color}          % Verwendung von Farbe 

\usepackage{acronym}        % für ein Abkürzungsverzeichnis

\usepackage[                % Biblatex
  backend=biber,
  bibstyle=lib/template/_dhbw_authoryear,maxbibnames=99,
  citestyle=authoryear,
  uniquename=true, useprefix=true,
  bibencoding=utf8]{biblatex}
%kein Punkt am Ende bei \footcite
%http://www.golatex.de/footcite-ohne-punkt-am-schluss-t4865.html
\renewcommand{\bibfootnotewrapper}[1]{\bibsentence#1}


%Reihenfolge der Autorennamen
%   
% http://golatex.de/viewtopic,p,80448.html#80448
% Argumente: siehe http://texwelt.de/blog/modifizieren-eines-biblatex-stils/
\DeclareNameFormat{sortname}{% Bibliographie
  \ifnum\value{uniquename}=0 % Normalfall
    \ifuseprefix%
      {%
         \usebibmacro{name:family-given}
           {\namepartfamily}
           {\namepartgiveni}
           {\namepartprefix}
           {\namepartsuffixi}%
       }
      {%
         \usebibmacro{name:family-given}
           {\namepartfamily}
           {\namepartgiveni}
           {\namepartprefixi}
           {\namepartsuffixi}%
       }%
  \fi
  \ifnum\value{uniquename}=1% falls nicht eindeutig, abgek. Vorname 
      {%
         \usebibmacro{name:family-given}
           {\namepartfamily}
           {\namepartgiveni}
           {\namepartprefix}
           {\namepartsuffix}%
       }%
  \fi
  \ifnum\value{uniquename}=2% falls nicht eindeutig, ganzer Vorname 
      {%
         \usebibmacro{name:family-given}
           {\namepartfamily}
           {\namepartgiven}
           {\namepartprefix}
           {\namepartsuffix}%
       }%
  \fi   
  \usebibmacro{name:andothers}}

\DeclareNameFormat{labelname}{% für Zitate
  \ifnum\value{uniquename}=0 % Normalfall
    \ifuseprefix%
      {%
         \usebibmacro{name:family-given}
           {\namepartfamily}
           {\empty}
           {\namepartprefix}
           {\namepartsuffixi}%
       }
      {%
         \usebibmacro{name:family-given}
           {\namepartfamily}
           {\empty}
           {\namepartprefixi}
           {\namepartsuffixi}%
       }%
  \fi
  \ifnum\value{uniquename}=1% falls nicht eindeutig, abgek. Vorname 
      {%
         \usebibmacro{name:family-given}
           {\namepartfamily}
           {\namepartgiveni}
           {\namepartprefix}
           {\namepartsuffix}%
       }%
  \fi
  \ifnum\value{uniquename}=2% falls nicht eindeutig, ganzer Vorname 
      {%
         \usebibmacro{name:family-given}
           {\namepartfamily}
           {\namepartgiven}
           {\namepartprefix}
           {\namepartsuffix}%
       }%
  \fi   
  \usebibmacro{name:andothers}}
      
  
\DeclareFieldFormat{extrayear}{% = the 'a' in 'Jones 1995a'
  \iffieldnums{labelyear}
    {\mknumalph{#1}}
    {\mknumalph{#1}}}        

\renewcommand*{\multinamedelim}{\addslash}
\renewcommand*{\finalnamedelim}{\addslash}
\renewcommand*{\multilistdelim}{\addslash}
\renewcommand*{\finallistdelim}{\addslash}

\renewcommand{\nameyeardelim}{~}

% Literaturverzeichnis: Doppelpunkt zwischen Name (Jahr): Rest 
% http://de.comp.text.tex.narkive.com/Tn1HUIXB/biblatex-authoryear-und-doppelpunkt
\renewcommand{\labelnamepunct}{\addcolon\addspace}

% damit die Darstellung für Vollzitate von Primärquellen in 
% Fußnoten später auf "nicht fett" geändert werden kann 
% (nur für Zitate von Sekundärliteratur relevant)
\newcommand{\textfett}[1]{\textbf{#1}}

% für Zitate von Sekundärliteratur:
\newcommand{\footcitePrimaerSekundaer}[4]{%
  \renewcommand{\textfett}[1]{##1}%
  \footnote{\fullcite[#2]{#1}, zitiert nach \cite[#4]{#3}}%  
  \renewcommand{\textfett}[1]{\textbf{##1}}%
}

% Im Literaturverzeichnis: Autor (Jahr) fett
\renewbibmacro*{author}{%
  \ifboolexpr{%
    test \ifuseauthor%
    and
    not test {\ifnameundef{author}}
  }
    {\usebibmacro{bbx:dashcheck}
       {\bibnamedash}
       {\usebibmacro{bbx:savehash}%
        \textfett{\printnames{author}}%
        \iffieldundef{authortype}
          {\setunit{\addspace}}
          {\setunit{\addcomma\space}}}%
     \iffieldundef{authortype}
       {}
       {\usebibmacro{authorstrg}%
        \setunit{\addspace}}}%
    {\global\undef\bbx@lasthash
     \usebibmacro{labeltitle}%
     \setunit*{\addspace}}%
  \textfett{\usebibmacro{date+extrayear}}}

% Sonderfall: Quelle ohne Autor, aber mit Herausgeber
% Name des Herausgebers wird fett gedruckt
\renewbibmacro*{bbx:editor}[1]{%
  \ifboolexpr{%
    test \ifuseeditor%
    and
    not test {\ifnameundef{editor}}
  }
    {\usebibmacro{bbx:dashcheck}
       {\bibnamedash}
       {\textfett{\printnames{editor}}%
        \setunit{\addcomma\space}%
        \usebibmacro{bbx:savehash}}%
     \usebibmacro{#1}%
     \clearname{editor}%
     \setunit{\addspace}}%
    {\global\undef\bbx@lasthash
     \usebibmacro{labeltitle}%
     \setunit*{\addspace}}%
  \textfett{\usebibmacro{date+extrayear}}}

% Anpassungen für deutsche Sprache
\DefineBibliographyStrings{ngerman}{%
	nodate = {{o.J.}},
	urlseen = {{Abruf:}},
	ibidem = {{ebenda}}
}

% keine Anführungszeichen beim Titel im Literaturverzeichnis
\DeclareFieldFormat[article,book,inbook,inproceedings,manual,misc,phdthesis,thesis,online,report]{title}{#1\isdot}

\newcommand{\literaturverzeichnis}{%
% nur Literaturverzeichnis
% (als eigenes Kapitel)
\phantomsection
\addcontentsline{toc}{chapter}{Literaturverzeichnis}
\spezialkopfzeile{Literaturverzeichnis}
\defbibheading{lit}{\chapter*{Literaturverzeichnis}}
\label{chapter:quellen}
\printbibliography[heading=lit,notkeyword=ausblenden]
} % mit DHBW-spezifischen Einstellungen

\usepackage{hyperref}       % URL-Formatierung, klickbare Verweise

\usepackage{tocloft}        % für Verzeichnis der Anhänge

\newcounter{anhcnt}
\setcounter{anhcnt}{0}
\newlistof{anhang}{app}{}

\newcommand{\anhang}[1]{%
  \refstepcounter{anhcnt}
  \setcounter{anhteilcnt}{0}
  \section*{Anhang \theanhcnt: #1}
  \addcontentsline{app}{section}{\protect\numberline{Anhang \theanhcnt}#1}\par
}

\newcounter{anhteilcnt}
\setcounter{anhteilcnt}{0}

\newcommand{\anhangteil}[1]{%
  \refstepcounter{anhteilcnt}
  \subsection*{Anhang~\arabic{anhcnt}/\arabic{anhteilcnt}: #1}
  \addcontentsline{app}{subsection}{\protect\numberline{Anhang \theanhcnt/\arabic{anhteilcnt}}#1}\par
}

\renewcommand{\theanhteilcnt}{Anhang \theanhcnt/\arabic{anhteilcnt}}

% vgl. S. 4 Paket-Beschreibung tocloft 	
% Einrückungen für Anhangverzeichnis
\makeatletter
\newcommand{\abstaendeanhangverzeichnis}{
  \renewcommand*{\l@section}{\@dottedtocline{1}{0em}{5.5em}}
  \renewcommand*{\l@subsection}{\@dottedtocline{2}{2.3em}{6.5em}}
}
\makeatother

% Abbildungs- und Tabellenverzeichnis
% Bezeichnungen
\renewcaptionname{ngerman}{\figurename}{Abb.}
\renewcaptionname{ngerman}{\tablename}{Tab.}
% Einrückungen
\makeatletter
\renewcommand*{\l@figure}{\@dottedtocline{1}{0em}{2.3em}}
\renewcommand*{\l@table}{\@dottedtocline{1}{0em}{2.3em}}
\makeatother


\usepackage{chngcntr}                % fortlaufende Zähler für Fußnoten, Abbildungen und Tabellen
\counterwithout{figure}{chapter}
\counterwithout{table}{chapter}
\counterwithout{footnote}{chapter}

\usepackage[automark]{scrlayer-scrpage}
%% Definitionen für Kopf- und Fußzeile auf normalen Seiten
\defpagestyle{kopfzeile}
{% Kopfdefinition
  (\textwidth,0pt)    % Länge der oberen Linie,Dicke der oberen Linie       
  {} % Definition für linke Seiten im doppelseitigen Layout
  {} % Definition für rechte Seiten im doppelseitigen Layout      
  {  % Definition für Seiten im einseitigen Layout
	\makebox[0pt][l]{\rightmark}% 
	\makebox[\linewidth]{}% 
  }        
  (\textwidth, 0.4pt) % Untere Linienlänge, Untere Liniendicke
}
{% Fußdefinition
  (\textwidth,0pt)    % Obere Linienlänge, Obere Liniendicke
  {} % Definition für linke Seiten im doppelseitigen Layout
  {} % Definition für rechte Seiten im doppelseitigen Layout
  {  % Definition für Seiten im einseitigen Layout
    \makebox[\linewidth]{}%
    \makebox[0pt][r]{\pagemark}%
  }
  (\textwidth, 0pt)   % Länge der unteren Linie,Dicke der unteren Linie
}

%% Definitionen für Kopf- und Fußzeile auf ersten Seiten eines Kapitels
\defpagestyle{kapitelkopfzeile}
{% Kopfdefinition
  (\textwidth,0pt)    % Länge der oberen Linie,Dicke der oberen Linie       
  {} % Definition für linke Seiten im doppelseitigen Layout
  {} % Definition für rechte Seiten im doppelseitigen Layout      
  {}  % Definition für Seiten im einseitigen Layout
  (\textwidth, 0pt) % Untere Linienlänge, Untere Liniendicke
}
{% Fußdefinition
  (\textwidth,0pt)    % Obere Linienlänge, Obere Liniendicke
  {} % Definition für linke Seiten im doppelseitigen Layout
  {} % Definition für rechte Seiten im doppelseitigen Layout
  {  % Definition für Seiten im einseitigen Layout
    \makebox[\linewidth]{}%
    \makebox[0pt][r]{\pagemark}%
  }
  (\textwidth, 0pt)   % Länge der unteren Linie,Dicke der unteren Linie
}

%% Definitionen für Kopf- und Fußzeile im Anhang und bei Quellenverzeichnisse
\newcommand{\spezialkopfzeileBezeichnung}{}
\defpagestyle{spezialkopfzeile}
{% Kopfdefinition
  (\textwidth,0pt)    % Länge der oberen Linie,Dicke der oberen Linie       
  {} % Definition für linke Seiten im doppelseitigen Layout
  {} % Definition für rechte Seiten im doppelseitigen Layout      
  {  % Definition für Seiten im einseitigen Layout
	\makebox[0pt][l]{\spezialkopfzeileBezeichnung}% 
	\makebox[\linewidth]{}% 
  }        
  (\textwidth, 0.4pt) % Untere Linienlänge, Untere Liniendicke
}
{% Fußdefinition
  (\textwidth,0pt)    % Obere Linienlänge, Obere Liniendicke
  {} % Definition für linke Seiten im doppelseitigen Layout
  {} % Definition für rechte Seiten im doppelseitigen Layout
  {  % Definition für Seiten im einseitigen Layout
    \makebox[\linewidth]{}%
    \makebox[0pt][r]{\pagemark}%
  }
  (\textwidth, 0pt)   % Länge der unteren Linie,Dicke der unteren Linie
}
            
\newcommand\spezialkopfzeile[1]{%
  \renewcommand\spezialkopfzeileBezeichnung{#1}
  \pagestyle{spezialkopfzeile}
}
                
% Standard-Pagestyle auswählen
\pagestyle{kopfzeile}

% keine Kopfzeile anzeigen auf Seiten, auf denen ein 
% Kapitel beginnt oder das Inhalts-/Abbildungs-/Tabellenverzeichnis steht 
\renewcommand{\chapterpagestyle}{kapitelkopfzeile}
\tocloftpagestyle{kapitelkopfzeile}

		 % für schöne Kopfzeilen 

\usepackage{textcomp}            % erlaubt EUR-Zeichen in Eingabedatei
\usepackage{eurosym}             % offizielles EUR-Symbol in Ausgabe
\renewcommand{\texteuro}{\euro}  % ACHTUNG: nach hyperref aufrufen!

\usepackage{scrhack}             % stellt Kompatibilität zw. KOMA-Script
% (scrreprt) und anderen Paketen her

% Anpassung der Abstände bei Kapitelüberschriften
% (betrifft auch Inhalts-, Abbildungs- und Tabellenverzeichnis)
\renewcommand*\chapterheadstartvskip{\vspace*{-\topskip}}
\newcommand{\myBeforeTitleSkip}{1mm}
\newcommand{\myAfterTitleSkip}{10mm}
\setlength\cftbeforetoctitleskip{\myBeforeTitleSkip}
\setlength\cftbeforeloftitleskip{\myBeforeTitleSkip}
\setlength\cftbeforelottitleskip{\myBeforeTitleSkip}

\setlength\cftaftertoctitleskip{\myAfterTitleSkip}
\setlength\cftafterloftitleskip{\myAfterTitleSkip}
\setlength\cftafterlottitleskip{\myAfterTitleSkip}
%%% Ende der Präambel %%%


%%% Name der eigenen Literatur-Datenbank %%%
\bibliography{lib/bibliography.bib}

\begin{document}

%%% Deckblatt einbinden %%% 
% Typ der Arbeit (für Deckblatt und ehrenwörtliche Erklärung) - bitte Zutreffendes auswählen
% 1. Projektarbeit
% 2. Projektarbeit
% Seminararbeit
% Bachelorarbeit
\newcommand{\bootstrapPaperType}{Benutzerdefinierter Typ}

% Thema der Arbeit (für ehrenwörtliche Erklärung, ohne Umbrüche)
\newcommand{\bootstrapPaperTitle}{Mein Titel}

% Ggf. Untertitel der Arbeit (falls dies nicht benötigt wird, das zweite Klammerpaar leer lassen, nicht das Kommando löschen!)
\newcommand{\bootstrapPaperSubtitle}{Untertitel}

% Vorname, Name der Autorin/des Autors (für Titelseite und Metadaten)
\newcommand{\bootstrapAuthor}{Vorname Nachname}

% Abgabedatum
\newcommand{\bootstrapDueDate}{\today}

% Fakultät der Arbeit festlegen
\newcommand{\bootstrapFaculty}{Wirtschaft}

% Studiengang festlegen.
\newcommand{\bootstrapMajor}{Wirtschaftsinformatik}

% Kurs festlegen
\newcommand{\bootstrapCourseID}{...}

% Daten bzgl der Betreuung
\newcommand{\bootstrapCompanyAdvisorGender}{f} % m/f/x für die verschiedenen Geschlechter
\newcommand{\bootstrapCompanyName}{Name des Unternehmens}
\newcommand{\bootstrapCompanyAdvisorDetails}{Titel, Vorname und Nachname der Betreuer*in}
\newcommand{\bootstrapUniversityAdvisorDetails}{Titel, Vorname und Nachname der Betreuer*in}
\newcommand{\bootstrapCompanyAdvisorPosition}{Funktion der Betreuerin/des Betreuers}
\newcommand{\bootstrapUniversityAdvisorPosition}{Funktion der Betreuerin/des Betreuers}


% Support command for generating the gender conform advisor term.
\newcommand{\genderedAdvisor}[1]{
    \ifthenelse{\equal{#1}{m}}
    {Betreuer}
    {\ifthenelse{\equal{#1}{f}}{Betreuerin}{Betreuungsperson}}
}

% Support command to generate the gender conform advisor term including an article.
\newcommand{\genderedAdvisorArticle}[1]{
    \ifthenelse{\equal{#1}{m}}
    {des Betreuers}
    {\ifthenelse{\equal{#1}{f}}{der Betreuerin}{der Betreuungsperson}}
}

\thispagestyle{empty}

\begin{spacing}{1}
    \begin{center}
        ~\vspace{0mm}

% Titel der Arbeit
        {\sffamily
        \LARGE
        \textbf{\bootstrapPaperTitle}

        \bigskip
        \textbf{\bootstrapPaperSubtitle}
        }

        \vspace{15mm}

% Typ der Arbeit
        {\Large \bootstrapPaperType}

        \vspace{1cm}

% Datum festlegen
        vorgelegt am\bootstrapDueDate

        \vspace{15mm}

% Fakultät festlegen
        Fakultät\bootstrapFaculty
        \medskip

% Studiengang festlegen
        Studiengang\bootstrapMajor
        \medskip

% Kurs festlegen
        Kurs\bootstrapClass

        \vspace{10mm}

        von

        \vspace{10mm}

% Autor/in/en/innen festlegen
        {\large\textsc{\bootstrapAuthor}}

        \vspace{10mm}
    \end{center}

    \vfill


% HIER EDITIEREN: Name des Unternehmens, Name der Betreuerin/des Betreuers
    \begin{tabular}{ll}
        \genderedSupport{\bootstrapCompanyAdvisorGender}in der Ausbildungsstätte: & DHBW Stuttgart:                     \\
        \hspace{0.4\linewidth}                                                    &                                     \\
        \bootstrapCompanyName                                                     & \bootstrapUniversityAdvisorDetails  \\
        \bootstrapCompanyAdvisorDetails                                           & \bootstrapUniversityAdvisorPosition \\
        \bootstrapCompanyAdvisorPosition \\
        Unterschrift\genderedSupportArticle{\bootstrapCompanyAdvisorGender} \\
    \end{tabular}


    \vspace{1cm}
%(etwas Platz für die Unterschrift der Betreuerin/des Betreuers aus der Ausbildungsstätte)
\end{spacing}

% falls ein Vertraulichkeitsvermerk erforderlich ist,
% die Kommentarzeichen in den nachfolgenden Zeilen entfernen:

%\begin{center}
%\small
%\textbf{Vertraulichkeitsvermerk}:
%Der Inhalt dieser Arbeit darf weder als Ganzes noch in Auszügen \\
%Personen außerhalb des Prüfungs- und Evaluationsverfahrens zugänglich gemacht werden, sofern keine anders lautende Genehmigung des Dualen Partners vorliegt. 
%\end{center}

% Meta-Daten für PDF-Datei basierend auf obigen Angaben
\hypersetup{pdftitle={\bootstrapPaperTitle}}
\hypersetup{pdfauthor={\bootstrapAuthor}}
\hypersetup{pdfsubject={\bootstrapPaperType\ DHBW Stuttgart \the\year}}


%%% Umstellung der Seiten-Nummerierung auf i, ii, iii ... %%%
\pagenumbering{Roman}

%%% Abstract einbinden (optionale Kurzfassung Ihrer Arbeit) %%%
% \begin{abstract}
\thispagestyle{kapitelkopfzeile}
\textbf{\LaTeX-Vorlage für Projekt-, Seminar- und Bachelorarbeiten}

Bei dem vorliegenden Dokument handelt es sich um eine Vorlage, die
für Projekt-, Seminar- und Bachelorarbeiten im Studiengang
Wirtschaftsinformatik der DHBW Stuttgart verwendet werden kann.

Sie setzt die technischen Vorgaben der Zitierrichtlinien\footnote{Sie finden diese unter \enquote{Prüfungsleistungen} im Studierendenportal (\url{https://studium.dhbw-stuttgart.de/winf/pruefungsleistungen/}).} des Studiengangs
(Stand: 01/2020) um.

\emph{Hinweise:} Bitte lesen Sie sich die Zitierrichtlinien unbedingt genau durch. Dieses Dokument ersetzt keine Anleitung oder Einführung in \LaTeX,
für die Nutzung sind daher gewisse Vorkenntnisse unerlässlich. Ein Einstieg in 
\LaTeX\ ist aber weniger schwierig, als es vielleicht auf den ersten Blick scheint
und lohnt sich für das Verfassen wissenschaftlicher Arbeiten in jedem Fall.\footnote{%
so auch \url{http://www.spiegel.de/netzwelt/tech/textsatz-keine-angst-vor-latex-a-549509.html}} 
Als Hilfestellung beim Schreiben eines Dokuments habe ich einen zweiseitigen kompakten \LaTeX-Spickzettel erstellt, der über Moodle verfügbar ist.

Ihre Rückmeldungen und Anregungen zu dieser Vorlage nehme ich gerne per E-Mail an die Adresse
\url{tobias.straub@dhbw-stuttgart.de} entgegen.

--- Prof. Dr. Tobias Straub

\vspace{5em}

\begin{center}\small
\begin{tabular}{ccl}
\multicolumn{3}{c}{\textbf{Versionshistorie}}\\
\hline
1.0	& 05.02.2015 & erste Fassung \\
\hline
1.1 & 16.02.2015 & siehe~\ref{anhang:ReleaseNotes11} \\
\hline
1.2 & 20.04.2015 & siehe~\ref{anhang:ReleaseNotes12} \\
\hline
1.3 & 20.02.2016 & siehe~\ref{anhang:ReleaseNotes13} \\
\hline
1.4 & 24.07.2017 & siehe~\ref{anhang:ReleaseNotes14} \\
\hline
1.5 & 07.01.2018 & siehe~\ref{anhang:ReleaseNotes15} \\
\hline
1.6 & 07.04.2018 & siehe~\ref{anhang:ReleaseNotes16} \\
\hline
1.7 & 12.02.2019 & siehe~\ref{anhang:ReleaseNotes17} \\
\hline
1.8 & 10.02.2020 & siehe~\ref{anhang:ReleaseNotes18} \\
\end{tabular}
\end{center}

\end{abstract}


% \cleardoublepage

\begin{spacing}{1}

  %%% Inhaltsverzeichnis %%%
  \tableofcontents
  \clearpage

  %%% Abkürzungsverzeichnis %%%
  \chapter*{Abkürzungsverzeichnis}
\addcontentsline{toc}{chapter}{Abkürzungsverzeichnis}

\begin{acronym}[DHBW] 
% Argument definiert die Breite der ersten Spalte anhand des längsten vorkommenden Eintrags
\acro{CRM}{Customer Relationship Management}
\acro{DHBW}{Duale Hochschule Baden-Württemberg}
\end{acronym}

  \clearpage

  \thispagestyle{kapitelkopfzeile}

  %%% Abbildungsverzeichnis %%%
  \listoffigures
  \phantomsection

  %%% Einfügen des Abbildungsverzeichnisses in das Inhaltsverzeichnis %%%
  \addcontentsline{toc}{chapter}{Abbildungsverzeichnis}
  \clearpage

  %%% Tabellenverzeichnis %%%
  \listoftables
  \phantomsection

  %%% Einfügen des Tabellenverzeichnisses in das Inhaltsverzeichnis %%%
  \addcontentsline{toc}{chapter}{Tabellenverzeichnis}
  \cleardoublepage

\end{spacing}

%%% Umstellung der Seiten-Nummerierung auf 1, 2, 3 ... %%%
\pagenumbering{arabic}

%%% Beginn des eigentlichen Inhalts %%%
% Empfehlung: strukturieren Sie Ihren Text in einzelnen Dateien und binden Sie diese hier mit   \input{main_part/dateiname.tex} ein

% Cheat Sheet mit einem Überblick über die wichtigsten Befehle
\chapter{Cheat Sheet}\label{chapter:cheat_sheet}

Hier findet Ihr einen Überblick über die wichtigsten \LaTeX\ Befehle.

\textbf{Fetter Text} \\
\textit{Kursiver Text} \\
\underline{Unterstrichener Text} \\

\section{Einfügen von Aufzählungen}\label{section:aufzaehlungen}

\begin{itemize}
    \item Erstes Element
    \item Zweites Element
    \item Drittes Element
\end{itemize}

\section{Einfügen von Bildern}\label{section:bilder}

% Einfügen einer Grafik
\begin{figure}[htb]
    % Zentrierung
    \centering
    % Einfügen der Datei, mit angepasster Höhe
    \includegraphics[height=5cm]{lib/graphics/dhbw.png}
    % Titel und Label der Grafik
    \caption[Logo der DHBW]{Logo der DHBW.\footnotemark}
    \label{abb:DHBWLogo}
\end{figure}
\footnotetext{Logo der DHBW}

\section{Einfügen von Tabellen}\label{section:tabellen}

% Einfügen einer Tabelle
\begin{table}[htb]
    % Zentrierung
    \centering
    % Beginn der eigentlichen Tabelle
    % Im zweiten Klammerpaar, Definition des Tabellendesign
    % Ausrichtung (l, c, r) und vertikale Trennung (|)
    \begin{tabular}{lcr}
        % Beginn einer neuen Spalte mit & und einer neuen Zeile mit \\
        links  & Mitte  & rechts \\
        % Einfügen einer horizontalen Linie
        \hline
        Muster & Muster & Muster \\
    \end{tabular}
    % Titel und Label der Tabelle
    \caption{Kleine Beispiel-Tabelle.}
    \label{tab:BeispielTabelleKlein}
\end{table}

\section{Einfügen von Source Code}\label{section:source_code}

% Definition, welche Programmiersprache verwendet wird
\lstset{language=Java}

In \LaTeX\ kann auch Source Code referenziert werden.

Dieser kann direkt in das Dokument geschrieben werden:

% Beginn des Code Blocks
\begin{lstlisting}
        if(1 > 0) {
        System.out.println("OK"); 
        } else {
        System.out.println("merkwuerdig");
        }
    \end{lstlisting}

\section{Einfügen von Abkürzungen}\label{section:abkuerzungen}

Beim ersten mal wird die Abkürzung ausgeschrieben: \ac{DHBW}

Bei allen weiteren Verwendungen wird nur die Abkürzung mit Link dargestellt: \ac{DHBW}

\section{Einfügen von Fußnoten und Zitaten}\label{section:zitate}

Fußnoten können auf verschiedene Weisen in \LaTeX\ eingefügt werden.

Hier wird nur eine beispielhafte Variante gezeigt!

\begin{itemize}
    \item 'Direktes Zitate'\footnote{\cite[S. 114ff.]{Mayring2002}}
    \item Indirektes Zitate\footnote{Vgl. \cite[S. 114ff.]{Mayring2002}}
    \item Sekundäres indirektes Zitat\footnote{Vgl. \cite[S. 114ff.]{Mayring2002} nach \cite{Endres}}
    \item Einzelnes sekundäres Zitat\footcite[Vgl.][S.1]{W3SchoolUnderscore}
\end{itemize}

\section{Einfügen von Referenzen}\label{section:referenzen}

Auf Inhalte kann referenziert werden. Hierzu müssen diese Inhalte nur mit einem Label versehen werden.

\begin{itemize}
    \item Bilder\ref{abb:DHBWLogo}
    \item Kapitel\ref{chapter:cheat_sheet}
    \item Abschnitte\ref{section:bilder}
\end{itemize}


\clearpage

% Beispieldateien
% \input{examples/einleitung.tex}
% \blinddocument
% \input{examples/text_mit_zitaten.tex}
% \chapter{Beispiele für Abbildungen und Tabellen}\label{chapter:abbildungenTabellen}

Hier finden Sie Beispiele für Abbildungen, Tabellen, Formelsatz und  Source Code.

\section{Abbildungen}
In diesem Abschnitt gibt die Abbildungen~\ref{abb:Logo2cmHoch} und~\ref{abb:Logo2cmBreit}, die beide das Logo der DHBW zeigen.

\begin{figure}[htb]
    \centering
    \includegraphics[height=2cm]{lib/graphics/dhbw.png}
    \caption[DHBW-Logo 2cm hoch]{DHBW-Logo 2cm hoch.\footnotemark}
    \label{abb:Logo2cmHoch}
\end{figure}
\footnotetext{Mit Änderungen entnommen aus: \cite{OhneAutorenOhneJahr}}

\lstset{language=TeX, % hervorzuhebende Keywords definieren
    morekeywords={footnotetext,footnotemark,footcite,caption}
}

\emph{Spezialfall:} Sofern \emph{innerhalb} der Bezeichnung einer Abbildung eine Fußnote angegeben oder eine Quelle referenziert werden soll, geschieht dies nicht per \lstinline|\footnote| oder \lstinline
|\footcite|. Vielmehr sind die Befehle \lstinline|\footnotemark| und \lstinline|\footnotetext| zu verwenden und außerdem das optionale Argument für \lstinline|\caption| anzugeben (vgl.\ Source Code).

\begin{figure}[htb]
    \centering
    \includegraphics[width=2cm]{lib/graphics/dhbw.png}
    \caption[DHBW-Logo 2cm breit.]{DHBW-Logo 2cm breit. (Quelle: DHBW\footnotemark)}
    \label{abb:Logo2cmBreit}
\end{figure}
\footnotetext{\url{www.dhbw.de}}



\section{Tabellen}

In diesem Abschnitt gibt es zwei Beispiel-Tabellen, nämlich auf Seite~\pageref{tab:BeispielTabelleKlein} und auf Seite~\pageref{tab:BeispielTabelleGroesser}.

\begin{table}[htb]
    \centering
    \begin{tabular}{lcr}
        links  & Mitte  & rechts \\
        \hline
        Muster & Muster & Muster \\
    \end{tabular}
    \caption{Kleine Beispiel-Tabelle.}
    \label{tab:BeispielTabelleKlein}
\end{table}

\begin{table}[htb]
    \centering
    \begin{tabular}{|l|l|c|l|r||l}
        \textbf{Spalte 1} & \textbf{Spalte 2} & \textbf{Spalte 3} & \textbf{Spalte 4} & \textbf{Spalte 5} & \textbf{Spalte 6} \\
        \hline
        a                 & b                 & c                 & d                 & e                 & f                 \\
        Test              & Test, Test        & Test, Test, Test  & ~                 & ~                 & ~                 \\
        1                 & 2                 & 3                 & 4                 & 5                 & 6                 \\
    \end{tabular}
    \caption{Größere Beispiel-Tabelle.}
    \label{tab:BeispielTabelleGroesser}
\end{table}

\section{Etwas Mathematik}

Eine abgesetzte Formel:
\[
    \int_a^b x^2 \: \mathrm{d} x = \frac{1}{3} (b^3 - a^3)
\]

Es ist $a^2+b^2 = c^2$ eine Formel im Text.

\section{Source Code}

Source Code-Blöcke können auf folgende Arten eingefügt werden:

\lstset{language=Java}

Direkt im \LaTeX-Source Code:
\begin{lstlisting}
if(1 > 0) {
  System.out.println("OK"); 
} else {
  System.out.println("merkwuerdig");
}
\end{lstlisting}

oder eingefügt aus einer externen Datei.
\lstinputlisting{includes/HelloWorld.java}

%%% Ende des eigentlichen Inhalts %%%

%%% Beginn des Anhangs %%%
\chapter*{Anhang}
\addcontentsline{toc}{chapter}{Anhang}

\lstset{language=TeX,
    morekeywords={anhang, anhangteil}
}

% Definition des Anhangverzeichnis
\section*{Anhangverzeichnis}
\vspace{-8em}
\abstaendeanhangverzeichnis
\listofanhang
\clearpage

% Konfiguration der speziellen Kopfzeile für den Anhang
\spezialkopfzeile{Anhang}

% Hauptteil des Anhangs
% Mit \anhang{Abschnitt des Anhangs} fügt man ein Kapitel in dem Anhang hinzu z. B. Interview Transkripte
\anhang{Interview Transkripte}

% Einfügen von weiteren Unterabschnitten, um den Anhang zu gliedern
% Mit \anhangteil{Name des Unterabschnitts} fügt man weitere Unterabschnitte hinzu
\anhangteil{Interview Transkript: Mitarbeiter eines Unternehmens}

% Es bietet sich an für bestimmte Dinge (wie die Namen von Interviewpartner) Variablen anzulegen
% \newcommand{\Interviewer}{\textbf{Interviewer: }}
% \newcommand{\Partner}{\textbf{Gesprächspartner: }}
\clearpage

%%% Quellenverzeichnisse %%%
\literaturverzeichnis
\cleardoublepage

%%% Erklärung %%%
\clearpage

\thispagestyle{empty}

{\LARGE\textsf{\textbf{Erklärung}}\bigskip}

% \typMeinerArbeit und \themaMeinerArbeit werden in src/config.tex definiert
Ich versichere hiermit, dass ich meine \bootstrapPaperType\ mit dem Thema: \emph{\bootstrapPaperTitle} selbstständig verfasst und keine anderen als die angegebenen Quellen und Hilfsmittel benutzt habe.
Ich versichere zudem, dass die eingereichte elektronische Fassung mit der gedruckten Fassung übereinstimmt.

\vspace{3cm}

\begin{center}
    \begin{tabular}{ccc}
        (Ort, Datum) & \hspace{0.3\linewidth} & (Unterschrift)
    \end{tabular}
\end{center}


\end{document}
